\documentclass[11pt]{letter}

\usepackage[margin=0.5in]{geometry}
\usepackage{hyperref}

\begin{document}

\begin{letter}{RBC Borealis}

\opening{Dear Hiring Team,}

I am writing to apply for the Machine Learning Researcher position at RBC Borealis. The role is exciting to me because it combines the two parts of machine learning work that I care about most: original research on difficult open problems and the translation of that research into systems that can create real product impact.

My M.Sc. research at the University of Calgary focused on machine learning and scalable data systems for large geospatial and Earth observation datasets. I independently formulated and built a DGGS-based framework for real-time multiresolution management of spatiotemporal satellite imagery, including C++/Python algorithms for encoding, storage, retrieval, temporal approximation, and performance evaluation. This work led to a first-author publication in Remote Sensing and gave me practical experience moving from a research question to a working software system, quantitative results, and a publishable technical contribution.

I also developed a deep learning framework for digital elevation model super-resolution. In that project, I trained PyTorch ResNet models, built data-processing pipelines around Google Earth Engine and Rasterio, and evaluated results using terrain-aware metrics such as RMSE, slope, TRI, TPI, and wavelet-based errors. A key part of the work was recognizing that random train-test splits produced overly optimistic results, then designing region-based validation to better measure generalization on real terrain. That experience maps closely to the kind of careful research judgment needed when working with real-world data rather than clean textbook benchmarks.

Beyond my graduate research, I bring a software engineering background that helps bridge research and implementation. I have built Django and PostgreSQL-backed APIs for NLP, speech, authentication, and payment workflows, and I have taught or supported courses involving deep learning, databases, large-scale data processing, and AWS-based data workflows. As a Neuromatch Academy Deep Learning TA, I mentored students through PyTorch, CNNs, attention and transformer models, diffusion models, and reinforcement learning projects.

I see a strong fit with RBC Borealis because your lab works at the intersection of publishable machine learning research, massive structured and unstructured datasets, and product-facing impact. My strongest current alignment is with Python, PyTorch, deep learning, scientific experimentation, research communication, and large-scale data processing. I have less direct experience in financial services and production transfer inside a bank, and I do not hold a PhD, but I do bring a demonstrated research track record, a first-author publication, and the ability to drive a technically complex research project independently.

I would be excited to contribute to RBC Borealis by developing rigorous, useful machine learning methods and helping move them toward impactful products. Thank you for considering my application.

\closing{Sincerely,

Amir Mirzai Golpayegani

\href{mailto:amirgolpaa24@gmail.com}{amirgolpaa24@gmail.com}}

\end{letter}

\end{document}
